\documentclass[xetex,mathserif,serif,handout]{beamer}

\usepackage{xunicode}
\usepackage{xltxtra}
\usepackage{color}
\usepackage{url}
\usepackage{listings}
\usepackage{fontspec}
\usepackage{geometry}
\usepackage{lastpage}
\usepackage{fancyhdr}
\usepackage{amsmath}
\usepackage{amsthm}
\usepackage{amssymb}
\usepackage{blkarray}
\usepackage{multicol}
\usepackage{relsize}

\definecolor{solarized@base03}{HTML}{002B36}
\definecolor{solarized@base02}{HTML}{073642}
\definecolor{solarized@base01}{HTML}{586e75}
\definecolor{solarized@base00}{HTML}{657b83}
\definecolor{solarized@base0}{HTML}{839496}
\definecolor{solarized@base1}{HTML}{93a1a1}
\definecolor{solarized@base2}{HTML}{EEE8D5}
\definecolor{solarized@base3}{HTML}{FDF6E3}
\definecolor{solarized@yellow}{HTML}{B58900}
\definecolor{solarized@orange}{HTML}{CB4B16}
\definecolor{solarized@red}{HTML}{DC322F}
\definecolor{solarized@magenta}{HTML}{D33682}
\definecolor{solarized@violet}{HTML}{6C71C4}
\definecolor{solarized@blue}{HTML}{268BD2}
\definecolor{solarized@cyan}{HTML}{2AA198}
\definecolor{solarized@green}{HTML}{859900}
\definecolor{yaleblue}{HTML}{0E4C92}

\setbeamertemplate{navigation symbols}{}
% \setbeamerfont{title}{family=\old}
% \setbeamerfont{author}{family=\tfont}%
% \setbeamerfont{frametitle}{family=\oldA}
% \setbeamerfont{date}{family=\dfont}

\setbeamertemplate{itemize items}{--}
\setbeamercolor*{item}{fg=black}

\defaultfontfeatures{Mapping=tex-text}
\hypersetup{pdfstartview={FitH}}

\newcommand{\old}[1]{\fontspec[Alternate=1,Ligatures={Common}]{Hoefler Text}\fontsize{18pt}{30pt}\selectfont #1}%
\newcommand{\oldA}[1]{\fontspec[Alternate=1,Ligatures={Common, Rare}]{Hoefler Text}\fontsize{12pt}{15pt}\selectfont #1}%
\newcommand{\oldB}[1]{\fontspec[Ligatures={Common}]{Didot}\fontsize{12pt}{15pt}\color{solarized@base02}\selectfont #1}%
\newcommand{\tfont}[1]{\fontspec[Alternate=1,Ligatures={Common}]{Hoefler Text}\fontsize{12pt}{20pt}\selectfont #1}%
\newcommand{\dfont}[1]{\fontspec[Ligatures={Common}]{Didot}\fontsize{12pt}{12pt}\selectfont #1}%

\newcommand{\minimize}{\mathop{\mathrm{minimize}}}
\newcommand{\argmin}{\mathop{\mathrm{arg\,min}}}
\newcommand{\argmax}{\mathop{\mathrm{arg\,max}}}
\newcommand{\st}{\mathop{\mathrm{subject\,\,to}}}

\newcommand\independent{\protect\mathpalette{\protect\independenT}{\perp}}
\def\independenT#1#2{\mathrel{\rlap{$#1#2$}\mkern2mu{#1#2}}}

\setlength{\parindent}{0pt}
\setlength{\parskip}{12pt}

\setromanfont [Ligatures={Common}, Numbers={OldStyle}, Variant=01,
 BoldFont={LinLibertine_RB.otf},
 ItalicFont={LinLibertine_RI.otf},
 BoldItalicFont={LinLibertine_RBI.otf}
 ]{LinLibertine_R.otf}



\begin{document}

%%%%%%%%%%%%%%%%%%%%%%%%%%%%%%%%%%%%%%%%%%%%%%%%%%%
\begin{frame}[fragile] \frametitle{}

\vfill

{\fontsize{0.7cm}{0cm}\selectfont Lecture 07 \\\vspace{0.2cm}
Hypothesis Testing with Multivariate Regression}\\\vspace{0.5cm}
23 September 2015

\vspace{2cm}

\begin{minipage}{0.6\textwidth}
Taylor B. Arnold \\
Yale Statistics \\
STAT 312/612
\end{minipage}
\hfill
\begin{minipage}{0.3\textwidth}\raggedleft
\includegraphics[scale=0.3]{../yale-logo.png}
\end{minipage}%

\end{frame}

%%%%%%%%%%%%%%%%%%%%%%%%%%%%%%%%%%%%%%%%%%%%%%%%%%%
\begin{frame}[fragile] \frametitle{}

{\color{yaleblue}\fontsize{16pt}{20pt}\selectfont Goals for today}

\begin{enumerate}
\item Review of assumptions and properties of linear model
\item The multivariate T-test
\item Hypothesis test simulations
\item The multivariate F-test
\end{enumerate}

\end{frame}

%%%%%%%%%%%%%%%%%%%%%%%%%%%%%%%%%%%%%%%%%%%%%%%%%%%
\begin{frame}[fragile] \frametitle{}

\begin{flushright}
{\color{yaleblue}\sc\fontsize{1cm}{0cm}\selectfont Review from Last Time}
\end{flushright}

\end{frame}


%%%%%%%%%%%%%%%%%%%%%%%%%%%%%%%%%%%%%%%%%%%%%%%%%%%
\begin{frame}[fragile] \frametitle{}

\begin{center}
{\Large Classical linear model assumptions}
\end{center}

{\bf I. Linearity} $\quad Y = X\beta + \epsilon$

{\bf II. Strict exogeneity} $\quad \mathbb{E} \left( \epsilon | X \right) = 0$

{\bf III. No multicollinearity} $\quad \mathbb{P} \left[ \text{rank} (X) = p \right] = 1$

{\bf IV. Spherical errors} $\quad \mathbb{V} \left( \epsilon | X \right) = \sigma^2 \mathbb{I}_n$

{\bf V. Normality} $\quad \epsilon | X \sim \mathcal{N} (0, \sigma^2 \mathbb{I}_n)$

\end{frame}


%%%%%%%%%%%%%%%%%%%%%%%%%%%%%%%%%%%%%%%%%%%%%%%%%%%
\begin{frame}[fragile] \frametitle{}

{\bf Finite sample properties}

Under assumptions I-III:

\hspace{1cm} (A) $\mathbb{E} (\widehat{\beta} | X) = \beta$

Under assumptions I-IV:

\hspace{1cm}  (B)  $\mathbb{V} (\widehat{\beta} | X) = \sigma^2 (X^t X)^{-1}$

\hspace{1cm}  (C) $\widehat{\beta}$ is the best linear unbiased estimator (Gauss-Markov)

\hspace{1cm}  (D) $Cov( \widehat{\beta}, r | X) = 0$

\hspace{1cm}  (E) $\mathbb{E} (s^2 | X) = \sigma^2$

Under assumptions I-V:

\hspace{1cm}  (F) $\widehat{\beta}$ achieves the Cramér–Rao lower bound

\end{frame}


%%%%%%%%%%%%%%%%%%%%%%%%%%%%%%%%%%%%%%%%%%%%%%%%%%%
\begin{frame}[fragile] \frametitle{}

\begin{flushright}
{\color{yaleblue}\sc\fontsize{1cm}{0cm}\selectfont The T-test}
\end{flushright}

\end{frame}

%%%%%%%%%%%%%%%%%%%%%%%%%%%%%%%%%%%%%%%%%%%%%%%%%%%
\begin{frame}[fragile] \frametitle{}

{\bf Hypothesis tests}

Consider testing the hypothesis $H_0: \beta_j = b_j$.

\pause Under assumptions I-V we have the following:
\begin{align*}
\left. \widehat{\beta}_j - b_j \right| X, H_0 \sim \mathcal{N} ( 0, \sigma^2 \left( (X^t X)^{-1}_{jj} \right))
\end{align*}

\end{frame}


%%%%%%%%%%%%%%%%%%%%%%%%%%%%%%%%%%%%%%%%%%%%%%%%%%%
\begin{frame}[fragile] \frametitle{}

{\bf Z-test}

This suggests the following test statistic:
\begin{align*}
z &= \frac{\widehat{\beta}_j - b_j}{\sqrt{\sigma^2  \left( (X^t X)^{-1}_{jj} \right)}}
\end{align*}
With,
\begin{align*}
\left. z \right| X, H_0 \sim \mathcal{N} (0, 1)
\end{align*}

\end{frame}

%%%%%%%%%%%%%%%%%%%%%%%%%%%%%%%%%%%%%%%%%%%%%%%%%%%
\begin{frame}[fragile] \frametitle{}

{\bf T-test}

As in the simple linear linear regression case, we generally
need to estimate $\sigma^2$ with $s^2$.

\pause This yields the following test statistic:
\begin{align*}
t &= \frac{\widehat{\beta}_j - b_j}{\sqrt{s^2  \left( (X^t X)^{-1}_{jj} \right)}} \\
&= \frac{\widehat{\beta}_j - b_j}{\text{S.E.}(\widehat{\beta}_j)}
\end{align*}


\end{frame}

%%%%%%%%%%%%%%%%%%%%%%%%%%%%%%%%%%%%%%%%%%%%%%%%%%%
\begin{frame}[fragile] \frametitle{}

{\bf T-test}

The test statistic has a $T$-distribution with $(n-p)$ degrees
of freedom under the null hypothesis:
\begin{align*}
\left. t \right| X, H_0 \sim t_{n-p}
\end{align*}
\pause This time around, we'll actually prove this.

\end{frame}

%%%%%%%%%%%%%%%%%%%%%%%%%%%%%%%%%%%%%%%%%%%%%%%%%%%
\begin{frame}[fragile] \frametitle{}

To start, we re-write the test statistic as:
\begin{eqnarray*}
t &=& \frac{\widehat{\beta} - b}{\sqrt{{\color{solarized@blue}\sigma^2}  \left( (X^t X)^{-1}_{jj} \right)}} \cdot
\sqrt{\frac{{\color{solarized@blue}\sigma^2}}{s^2}} \\ \pause
&=& \frac{z}{\sqrt{s^2 / \sigma^2}} \\ \pause
&=& \frac{z}{\sqrt{{\color{solarized@magenta}q} / (n-p)}}
\end{eqnarray*}
\pause Where ${\color{solarized@magenta}q} = r^t r / \sigma^2$.

\pause We need to show that (1) ${\color{solarized@magenta}q} | X \sim \chi^2_{n-p}$ and
(2) $z \independent {\color{solarized@magenta}q} | X$.

\end{frame}

%%%%%%%%%%%%%%%%%%%%%%%%%%%%%%%%%%%%%%%%%%%%%%%%%%%
\begin{frame}[fragile] \frametitle{}

{\bf Lemma} If $B$ is a symmetric idempotent matrix and
$u \sim \mathcal{N} (0, \mathbb{I}_n)$, then
$u^t B u \sim \chi^2_{\text{tr(B)}}$.

\end{frame}

%%%%%%%%%%%%%%%%%%%%%%%%%%%%%%%%%%%%%%%%%%%%%%%%%%%
\begin{frame}[fragile] \frametitle{}

{\it Proof:} The symmetric matrix $B$ can be written as
by its eigen-decompostion; for some orthonormal $Q$ matrix
and diagonal matrix $\Lambda$:
\begin{align*}
B &= Q^t \Lambda Q
\end{align*}

\end{frame}

%%%%%%%%%%%%%%%%%%%%%%%%%%%%%%%%%%%%%%%%%%%%%%%%%%%
\begin{frame}[fragile] \frametitle{}

Notice that as $B$ is idempotent:
\begin{eqnarray*}
(Q^t \Lambda Q) &=& (Q^t \Lambda Q) (Q^t \Lambda Q) \\ \pause
&=& Q^t \Lambda Q Q^t \Lambda Q \\ \pause
&=& Q^t \Lambda^2 Q \\
\end{eqnarray*}
\pause Since $Q^t = Q^{-1}$:
\begin{align*}
\Lambda &= \Lambda^2 \\
\end{align*}
\pause Therefore all the elements of $\Lambda$ are $0$ or $1$.

\end{frame}

%%%%%%%%%%%%%%%%%%%%%%%%%%%%%%%%%%%%%%%%%%%%%%%%%%%
\begin{frame}[fragile] \frametitle{}

By rotating the columns of $Q$, we can actually write:
\begin{align*}
\Lambda &= \left( \begin{array}{cc} \mathbb{I}_{\text{rank}(B)} & 0 \\ 0 & 0 \end{array} \right)
\end{align*}
\pause Notice that this also shows that:
\begin{align*}
\text{tr}(B) &= \text{tr}(Q^t \Lambda Q) \\
&= \text{tr}(\Lambda Q Q^t ) \\
&= \text{tr}(\Lambda) \\
&= \text{rank}(B)
\end{align*}

\end{frame}

%%%%%%%%%%%%%%%%%%%%%%%%%%%%%%%%%%%%%%%%%%%%%%%%%%%
\begin{frame}[fragile] \frametitle{}

Now let $v = Q u$. We can see that the mean of $v$ is zero:
\begin{align*}
\mathbb{E} v &= \mathbb{E} Q u \\
&= Q \mathbb{E} u \\
&= 0
\end{align*}
\pause And the variance of $v$ is:
\begin{align*}
\mathbb{E} v v^t &= Q \mathbb{E} (u u^t) Q \\
&= Q \mathbb{I}_n Q^t \\
&= \mathbb{I}_n
\end{align*}

\end{frame}


%%%%%%%%%%%%%%%%%%%%%%%%%%%%%%%%%%%%%%%%%%%%%%%%%%%
\begin{frame}[fragile] \frametitle{}

Now we calculate the quadtratic form $u^t B u$ by the
transform of $v$. Substitute that $u = (Q)^{-1}v = Q^tv$:
\begin{eqnarray*}
u^t B u &=& v^t Q B Q^t v \\ \pause
&=& v^t Q {\color{solarized@blue} (Q^t \Lambda Q)} Q^t v \\ \pause
&=& v^t {\color{solarized@magenta} Q Q^t} \Lambda {\color{solarized@magenta} Q Q^t} v \\ \pause
&=& v^t \Lambda v \\ \pause
&=& v^t \left( \begin{array}{cc} \mathbb{I}_{\text{tr}(B)} & 0 \\ 0 & 0 \end{array} \right) v \\ \pause
&=& \sum_{i=1}^{\text{tr}(B)} v_i^2 \pause \\
&\sim& \chi^2_{\text{tr}(B)}
\end{eqnarray*}
Which completes the lemma.

\end{frame}

%%%%%%%%%%%%%%%%%%%%%%%%%%%%%%%%%%%%%%%%%%%%%%%%%%%
\begin{frame}[fragile] \frametitle{}

(1) We know that $r^t r = \epsilon^t M \epsilon$, so:
\begin{eqnarray*}
q &=& \frac{r^t r}{\sigma^2} \\ \pause
&=& \frac{\epsilon^t}{\sigma} \cdot M \cdot \frac{\epsilon}{\sigma}
\end{eqnarray*}
\pause Under Assumption V,
$\frac{\epsilon}{\sigma} \sim \mathcal{N} (0, \mathbb{I}_n)$.
Therefore from the lemma, then $q \sim \chi^2_{\text{tr(M)}}$; last
time we showed the $\text{tr(M)} = n - p$, which
finishes the first part of the proof.

\end{frame}

%%%%%%%%%%%%%%%%%%%%%%%%%%%%%%%%%%%%%%%%%%%%%%%%%%%
\begin{frame}[fragile] \frametitle{}

(2) The random variables $\widehat{\beta}$ and $r$ are both
linear combinations of $\epsilon$, and are therefore jointly
normally distributed. As they are uncorrelated (problem set 2),
this implies that they are independent. Finally, this implies that
$z = f(\widehat{\beta})$ and $q = g(r)$ and themselves independent.

\end{frame}

%%%%%%%%%%%%%%%%%%%%%%%%%%%%%%%%%%%%%%%%%%%%%%%%%%%
\begin{frame}[fragile] \frametitle{}

So, therefore, we have:
\begin{align*}
t &= \frac{\widehat{\beta} - b}{\sqrt{{\color{solarized@blue}\sigma^2}  \left( (X^t X)^{-1}_{jj} \right)}} \cdot
      \sqrt{\frac{{\color{solarized@blue}\sigma^2}}{s^2}} \\
&= \frac{z}{{\color{solarized@magenta}q} / (n-p)} \\
&\sim t_{n-p}
\end{align*}

\end{frame}

%%%%%%%%%%%%%%%%%%%%%%%%%%%%%%%%%%%%%%%%%%%%%%%%%%%
\begin{frame}[fragile] \frametitle{}

\begin{flushright}
{\color{yaleblue}\sc\fontsize{1cm}{0cm}\selectfont Simulation}
\end{flushright}

\end{frame}

%%%%%%%%%%%%%%%%%%%%%%%%%%%%%%%%%%%%%%%%%%%%%%%%%%%
\begin{frame}[fragile] \frametitle{}

\begin{flushright}
{\color{yaleblue}\sc\fontsize{1cm}{0cm}\selectfont F-Test}
\end{flushright}

\end{frame}

%%%%%%%%%%%%%%%%%%%%%%%%%%%%%%%%%%%%%%%%%%%%%%%%%%%
\begin{frame}[fragile] \frametitle{}

{\bf F-test}

Now consider the hypothesis test $H_0: D\beta = d$ for a
matrix $D$ with $k$ rows and rank $k$.

\end{frame}

%%%%%%%%%%%%%%%%%%%%%%%%%%%%%%%%%%%%%%%%%%%%%%%%%%%
\begin{frame}[fragile] \frametitle{}

We'll form the following test statistic:
\begin{align*}
F &= \frac{(D\widehat{\beta} - d)^t [D (X^t X)^{-1} D^t ]^{-1} (D\widehat{\beta} - d) / k}{s^2}
\end{align*}
And prove that is has an F-distribution with $k$ and $n-p$ degrees of
freedom.

\end{frame}

%%%%%%%%%%%%%%%%%%%%%%%%%%%%%%%%%%%%%%%%%%%%%%%%%%%
\begin{frame}[fragile] \frametitle{}

We can re-write the test statistics as:
\begin{align*}
F &= \frac{{\color{solarized@cyan} w} / k}{{\color{solarized@magenta} q} / (n - p)}
\end{align*}
Where:
\begin{align*}
{\color{solarized@cyan} w} &= (D\widehat{\beta} - d)^t [\sigma^2 D (X^t X)^{-1} D^t ]^{-1} (D\widehat{\beta} - d) \\
{\color{solarized@magenta} q} &= r^t r / \sigma^2
\end{align*}

\pause
We've already shown that ${\color{solarized@magenta} q} | X \sim \chi^2_{n-p}$, and can see
that ${\color{solarized@magenta} q} \independent {\color{solarized@cyan} w}$ by the same argument as before. All that
is left is to show that ${\color{solarized@cyan} w} \sim \chi^2_{k}$.

\end{frame}

%%%%%%%%%%%%%%%%%%%%%%%%%%%%%%%%%%%%%%%%%%%%%%%%%%%
\begin{frame}[fragile] \frametitle{}

Let $v = D\widehat{\beta} - d$. Under the null hypothesis $v = D (\widehat{\beta} - \beta)$.
\pause Therefore:
\begin{eqnarray*}
\mathbb{V} (v | X) &=& \mathbb{V} (D (\widehat{\beta} - \beta) | X) \\ \pause
&=& D \mathbb{V} (\widehat{\beta} - \beta | X) D^t \\ \pause
&=& \sigma^2 D (X^t X)^{-1} D^t
\end{eqnarray*}

\pause
Now, for the multivariate normally distributed $v$ with zero mean,
we have:
\begin{align*}
w &= v^t \mathbb{V} (v | X)^{-1} v
\end{align*}
\end{frame}

%%%%%%%%%%%%%%%%%%%%%%%%%%%%%%%%%%%%%%%%%%%%%%%%%%%
\begin{frame}[fragile] \frametitle{}

Now, decompose $ \mathbb{V} (v | X)^{-1} = \Sigma^{-1}$
as $Q^t \Lambda Q$. Notice that this implies:
\begin{eqnarray*}
\Sigma^{-1} &=& Q^t \Lambda^{1/2} \Lambda^{1/2} Q \\ \pause
Q \Sigma^{-1} &=& \Lambda^{1/2} \Lambda^{1/2} Q  \\ \pause
\Lambda^{-1/2} Q \Sigma^{-1} &=& \Lambda^{1/2} Q
\end{eqnarray*}

\end{frame}

%%%%%%%%%%%%%%%%%%%%%%%%%%%%%%%%%%%%%%%%%%%%%%%%%%%
\begin{frame}[fragile] \frametitle{}

From here, we can write $w$ as the inner product of
a suitably defined vector $u$:
\begin{eqnarray*}
w &=& v^t \Sigma^{-1} v \\ \pause
&=& v^t Q^t \Lambda^{1/2} \Lambda^{1/2} Q v \\ \pause
&=& u^t u
\end{eqnarray*}
\pause With a mean of zero:
\begin{align*}
\mathbb{E} (u | X) &= \mathbb{E} (\Lambda^{1/2} Q v | X) \\
&= \Lambda^{1/2} Q \mathbb{E} (v | X) \\
&= 0
\end{align*}

\end{frame}

%%%%%%%%%%%%%%%%%%%%%%%%%%%%%%%%%%%%%%%%%%%%%%%%%%%
\begin{frame}[fragile] \frametitle{}

And variance of
\begin{eqnarray*}
\mathbb{E} (u u^t | X) &=& \mathbb{E} (\Lambda^{1/2} Q v v^t Q^t \Lambda^{1/2} | X) \\ \pause
&=& \Lambda^{1/2} Q \mathbb{E} (v v^t | X) Q^t \Lambda^{1/2} \\ \pause
&=& {\color{solarized@magenta} \Lambda^{1/2} Q} \Sigma Q^t \Lambda^{1/2} \\ \pause
&=& \Lambda^{-1/2} Q \Sigma^{-1}\Sigma Q^t \Lambda^{1/2} \\ \pause
&=& \Lambda^{-1/2} \Lambda^{1/2} \\ \pause
&=& \mathbb{I}_k
\end{eqnarray*}
And, finally:
\begin{align*}
w &= u u^t \sim \chi^2_{k}
\end{align*}
Which finishes the proof.

\end{frame}


%%%%%%%%%%%%%%%%%%%%%%%%%%%%%%%%%%%%%%%%%%%%%%%%%%%
\begin{frame}[fragile] \frametitle{}

{\bf An alternative F-test}

Now, consider the following estimator:
\begin{align*}
\widetilde{\beta} &= \argmin_b || y - Xb ||_2^2 , \quad \text{s.t.} \quad Db = d
\end{align*}
\pause We then define the restricted residuals
$\widetilde{r} = y - X\tilde{\beta}$.

\end{frame}

%%%%%%%%%%%%%%%%%%%%%%%%%%%%%%%%%%%%%%%%%%%%%%%%%%%
\begin{frame}[fragile] \frametitle{}

{\bf An alternative F-test}

An alternative expression for the $F$-test statistic is:
\begin{align*}
F &= \frac{(\tilde{r}^t \tilde{r} -  r^t r) / k }{r^t r / (n - p)}
\end{align*}
Conceptually, it should make sense that this is large whenever
the null hypothesis is false.

\end{frame}

%%%%%%%%%%%%%%%%%%%%%%%%%%%%%%%%%%%%%%%%%%%%%%%%%%%
\begin{frame}[fragile] \frametitle{}

{\bf An alternative F-test}

If we let $\text{SSR}_U$ be the sum of squared residuals of the unrestricted
model ($r^t r$) and $\text{SSR}_R$ be the sum of squared residuals of the
restricted model, then this can be re-written as:
\begin{align*}
F &= \frac{(\text{SSR}_R -  \text{SSR}_U) / k }{\text{SSR}_U / (n - p)}
\end{align*}
This is the way it is written in the homework and in the Fumio Hayashi
text. It is left for you to prove that this is equivalent to the
other F test.

\end{frame}

%%%%%%%%%%%%%%%%%%%%%%%%%%%%%%%%%%%%%%%%%%%%%%%%%%%
\begin{frame}[fragile] \frametitle{}

\begin{flushright}
{\color{yaleblue}\sc\fontsize{1cm}{0cm}\selectfont Application}
\end{flushright}

\end{frame}

\end{document}











